%% Figure: slope chart of detection rates on mined issue reproducers across versions.
\begin{figure}[t]
\centering
\begin{tikzpicture}
\begin{axis}[
    tccaxis, width=0.8\linewidth, height=4.0cm,
    xmin=-0.75, xmax=3.7, ymin=0, ymax=60,
    xtick={1,2}, xticklabels={PyTorch 2.10, PyTorch 2.14},
    ylabel={issues detected (\%)},
    clip=false,
]
  \addplot[tolorange, very thick, mark=*, mark size=2.2pt] coordinates {(1,45.7) (2,41.4)};
  \addplot[tolblue, very thick, mark=square*, mark size=2.2pt] coordinates {(1,40.7) (2,13.0)};
  \node[anchor=east, font=\scriptsize, tolorange!80!black] at (axis cs:0.9,48.5) {open issues: 32/70};
  \node[anchor=east, font=\scriptsize, tolblue!80!black] at (axis cs:0.9,37.5) {closed issues: 50/123};
  \node[anchor=west, font=\scriptsize, tolorange!80!black] at (axis cs:2.07,41.4) {29/70};
  \node[anchor=west, font=\scriptsize, tolblue!80!black] at (axis cs:2.07,13.0) {16/123};
  \draw[black!55, decorate, decoration={brace, amplitude=3pt, mirror}] (axis cs:2.55,14) -- (axis cs:2.55,40);
  \node[anchor=west, font=\scriptsize, black!75, align=left] at (axis cs:2.66,27) {35 verified\\buggy/fixed pairs};
\end{axis}
\end{tikzpicture}
\caption{RQ1 on 193 reproducers mined from the PyTorch issue tracker, executed on the same Linux CPU image with
two PyTorch versions. Open issues remain detectable; closed issues stop being detected once their fix ships.}
\label{fig:rq1-history}
\end{figure}
